\fancyhead[L]{\small FUNDAMENTAL EQUATIONS.}
\SourceLeaf{501}\pdfbookmark[0]{Fundamental Equations for the Motion of Heavy Bodies on the Rotating Earth.}{rotating-earth-title}
\begin{center}
{\large\bfseries FUNDAMENTAL EQUATIONS\par\medskip FOR THE MOTION OF HEAVY BODIES\par\medskip ON THE ROTATING EARTH.\par}
\bigskip
\textsc{Benzenberg}. Versuche über das Gesetz des Falls. 1804.
\par\bigskip
{\large\itshape Letter from Gauss to Benzenberg.\par}
\end{center}
\begin{flushright}Braunschweig, 1803. 2 February.\end{flushright}
—\quad—\quad—\quad There is an assumption in our friend \textsc{Olbers}'s theory that seems inadmissible to me. Namely, \textit{that the body remains in one plane during its fall.} In my opinion, one may \textit{not} assume this when taking air resistance into account, which is \textit{essential} here because the inferred southward deflection depends entirely on it. A simple consideration shows that the plane $(\mathrm{A})$ in which the body first begins to move passes through the centre of the Earth (or, more generally, the centre of attraction) and is perpendicular to the plane $(\mathrm{B})$ in which the meridian of the observation site lay at the start of the fall. But one can easily see that the particles of air pass obliquely through the plane $\mathrm{A}$ at every point except along the straight line where $\mathrm{A}$ is intersected by $\mathrm{B}$. The air therefore opposes the body's motion in a direction outside this plane $\mathrm{A}$ and drives it out of the plane to the north. It seemed to me that the effect would be just large enough to cancel the southward deflection inferred from the delay in the fall.

Your last letter prompted me to think about this again. In a spare half-hour I looked at the matter in quite a different way and derived the analytical equations,
\SourceLeaf{502}which describe the body's motion relative to the moving surface of the Earth, from the first principles of dynamics. To my surprise I found

1) the southward deflection again $0$ or quite imperceptible.

2) the eastward deflection only $\tfrac23$ of what Dr.\ \textsc{Olbers} found. In Dr.\ \textsc{Olbers}'s notation, if air resistance is neglected, it is
\[
=\frac{\tfrac43\pi\cos\psi.at}{86164}
\]
or, if air resistance is taken into account, to an approximation sufficient here,
\[
=\frac{\tfrac43\pi\cos\psi.t}{86164}(\tfrac32a'-\tfrac12a)
\]
where $a$ is the height through which the body would fall in time $t$ in a vacuum, and thus $=\tfrac12g'tt$ [where $\psi$ denotes the polar altitude of the observation site and $a'$ the actual height of fall].

On this basis, for your experiments, taking the length of the pendulum for Hamburg $=440.75$ lines (which makes $g'$ come out almost as Dr.\ \textsc{Olbers} takes it), I find \textit{the eastward deflection} to be $3.951$ \textit{Paris lines}. This agreed very closely with your experiments, — whereas the southward deflection does not agree with my results.

\begin{center}\resizebox{!}{58mm}{\begin{tikzpicture}[x=1bp,y=-1bp,line cap=round,line join=round]
\path[use as bounding box] (0,0) rectangle (395,752);
\begin{pgfinterruptboundingbox}
\draw[line width=2.3bp] (75,215) .. controls (140,109) and (253,62) .. (326,56);
\draw[line width=1.8bp] (75,215) -- (327,218);
\draw[line width=1.7bp] (327,56) -- (327,85);
\draw[line width=1.7bp] (327,106) -- (326,257);
\draw[line width=1.7bp] (326,270) -- (324,316);
\draw[line width=1.7bp] (324,326) -- (323,411);
\draw[line width=1.7bp] (323,438) -- (317,684);
\node[inner sep=0pt,outer sep=0pt] at (315,37) {\scalebox{3.8}{\ensuremath{m}}};
\node[inner sep=0pt,outer sep=0pt] at (45,201) {\scalebox{3.8}{\ensuremath{M}}};
\node[inner sep=0pt,outer sep=0pt] at (347,207) {\scalebox{3.8}{\ensuremath{R}}};
\node[inner sep=0pt,outer sep=0pt] at (316,708) {\scalebox{3.8}{\ensuremath{C}}};
\end{pgfinterruptboundingbox}
\end{tikzpicture}
}\end{center}
This difference in the eastward deflection — led me to examine Dr.\ \textsc{Olbers}'s conclusions about it more closely and to trace its cause. As it seems to me, the cause is that Dr.\ \textsc{Olbers} derives the body's actual eastward motion \textit{solely} from its initial tangential velocity, and subtracts from the resulting motion the simultaneous motion of the foot of the tower to obtain the apparent eastward motion. — But if the surface of the paper represents the plane $\mathrm{A}$ above, $\mathrm{C}$ the centre of the Earth, and $\mathrm{mM}$ the body's actual motion, then in my opinion we must bear in mind that the attraction towards $\mathrm{C}$ itself is not parallel to $\mathrm{mC}$ during the motion. Just for that rea\SourceLeaf{503}son the eastward velocity really decreases, so that when the body reaches $\mathrm{M}$ it has not gone as far east as it would have at its initial velocity. A calculation of this also shows me that taking this into account does indeed reduce the apparent eastward motion by one third.

\Needspace{7\baselineskip}
\begin{center}{\large\itshape Letter from Gauss to Benzenberg.}\end{center}
\begin{flushright}Braunschweig, 1803. 8 March.\end{flushright}
—\quad—\quad—\quad—\quad Eight days ago I sent our friend \textsc{Olbers} a short paper on the deflection of falling bodies. Today I receive his reply:

1) the eastward deflection is only $\tfrac23$ of the one he had calculated.

2) he has no objection to my conclusions that the southward deflection is $=0$. But he would like to know \textit{where} exactly \textit{his} reasoning goes wrong.

I would add the following.

Assuming my conclusions on the southward deflection are certain, it still seems to me that the deflection Dr.\ \textsc{Olbers} obtained comes from his assumption that, even with air resistance, the body \textit{remains} in the plane perpendicular to the meridian and passing through the centre of the Earth. It seems to me that this assumption needs to be justified, but I doubt it can be. As soon as the Earth has moved out of its initial position, the conical motion of the air makes the particles of air surrounding the body cross that plane at an angle that cannot be neglected. As a result, the body, whose motion the air does not oppose along that plane, leaves the plane towards the north. I still think this completely offsets the deflection due to the delay. It also seems likely to me that this is just what \textsc{Guglielmini} meant, and that he failed to win \textsc{Olbers}'s agreement only because he did not express himself precisely enough. I am confident, however, that either I shall come entirely round to Dr.\ \textsc{Olbers}'s view or Dr.\ \textsc{Olbers} will come entirely round to mine. —\quad—\quad—

\SourceLeaf{504}\Needspace{8\baselineskip}
\begin{center}{\large\itshape Fundamental Equations for the Motion of Heavy Bodies\par on the Rotating Earth.}\end{center}
The position of a point is determined in two ways.

\textit{First}, by its perpendicular distances $X,Y,Z$ from three mutually perpendicular \textit{fixed} planes. We place the common point of intersection of these planes, $C$, at any point on the Earth's axis. We set the plane for $Z$ parallel to the equator, the plane for $Y$ in the meridian that contains the body's initial position, and finally the plane for $X$ in the meridian perpendicular to the preceding one. The $Z$ distances are positive on the north side. The $X$ distances are positive on the side containing the body's initial position, and the $Y$ distances on the side towards which the rotation carries that initial position.

\textit{Second}, by its perpendicular distances $x,y,z$ from three mutually perpendicular \textit{moving} planes, that is, planes at rest relative to the Earth and rotating with it. It is most convenient to place their common point of intersection at the body's initial position. We set the plane for $z$ perpendicular to the apparent direction of gravity and the plane for $y$ in the meridian. This determines the plane for $x$, which is perpendicular to both. The poles of these three planes are thus the apparent zenith, the east point and the south point, respectively. \textit{These} poles also designate the sides of the planes on which the distances $z,y,x$ are taken as positive.

Now let $C$ have $x=a$, $(y=0)$, $z=-c$. Let the (apparent, northern) polar altitude of the observation site be $\varphi$ and the angle through which the Earth has moved eastward after time $t$ be $\theta$. Under these assumptions, the following equations are readily obtained.
\[
\left.\begin{aligned}
x&=X\sin\varphi\cos\theta+Y\sin\varphi\sin\theta-Z\cos\varphi+a\\
y&=-X\sin\theta+Y\cos\theta\\
z&=X\cos\varphi\cos\theta+Y\cos\varphi\sin\theta+Z\sin\varphi-c
\end{aligned}\right\}\quad.\quad.\quad.\quad.\quad.\quad[1]
\]
\[
\left.\begin{aligned}
X&=(x-a)\sin\varphi\cos\theta-y\sin\theta+(z+c)\cos\varphi\cos\theta\\
Y&=(x-a)\sin\varphi\sin\theta+y\cos\theta+(z+c)\cos\varphi\sin\theta\\
Z&=-(x-a)\cos\varphi+(z+c)\sin\varphi
\end{aligned}\right\}\quad.\quad.\quad.\quad[2]
\]

\SourceLeaf{505}The coordinates $X,Y,Z$ can be regarded either as functions of $t$ alone or as functions of the four variables $\theta,x,y,z$. Viewed in the second way, they therefore have four partial differentials. Thus
\[
\begin{aligned}
\mathrm{d}X&=(\tfrac{\mathrm{d}X}{\mathrm{d}t})\mathrm{d}t=(\tfrac{\mathrm{d}X}{\mathrm{d}\theta})\mathrm{d}\theta+(\tfrac{\mathrm{d}X}{\mathrm{d}x})\mathrm{d}x+(\tfrac{\mathrm{d}X}{\mathrm{d}y})\mathrm{d}y+(\tfrac{\mathrm{d}X}{\mathrm{d}z})\mathrm{d}z\\
\mathrm{d}Y&=\mathrm{etc.}
\end{aligned}
\]
The body's velocity, like its motion, can be resolved into three component velocities perpendicular to the planes for $X,Y,Z$. These components are therefore $(\tfrac{\mathrm{d}X}{\mathrm{d}t})$, $(\tfrac{\mathrm{d}Y}{\mathrm{d}t})$, $(\tfrac{\mathrm{d}Z}{\mathrm{d}t})$. The corresponding velocities of the element of air in which the body happens to be, with respect to the same planes, are evidently $(\tfrac{\mathrm{d}X}{\mathrm{d}\theta})\tfrac{\mathrm{d}\theta}{\mathrm{d}t}$, $(\tfrac{\mathrm{d}Y}{\mathrm{d}\theta})\tfrac{\mathrm{d}\theta}{\mathrm{d}t}$, $(\tfrac{\mathrm{d}Z}{\mathrm{d}\theta})\tfrac{\mathrm{d}\theta}{\mathrm{d}t}$. The body's \textit{relative} velocities in these three directions are consequently
\[
\begin{aligned}
(\tfrac{\mathrm{d}X}{\mathrm{d}x})\tfrac{\mathrm{d}x}{\mathrm{d}t}+(\tfrac{\mathrm{d}X}{\mathrm{d}y})\tfrac{\mathrm{d}y}{\mathrm{d}t}+(\tfrac{\mathrm{d}X}{\mathrm{d}z})\tfrac{\mathrm{d}z}{\mathrm{d}t}&=\xi=\sin\varphi\cos\theta\tfrac{\mathrm{d}x}{\mathrm{d}t}-\sin\theta\tfrac{\mathrm{d}y}{\mathrm{d}t}+\cos\varphi\cos\theta\tfrac{\mathrm{d}z}{\mathrm{d}t}\\
(\tfrac{\mathrm{d}Y}{\mathrm{d}x})\tfrac{\mathrm{d}x}{\mathrm{d}t}+(\tfrac{\mathrm{d}Y}{\mathrm{d}y})\tfrac{\mathrm{d}y}{\mathrm{d}t}+(\tfrac{\mathrm{d}Y}{\mathrm{d}z})\tfrac{\mathrm{d}z}{\mathrm{d}t}&=\eta=\sin\varphi\sin\theta\tfrac{\mathrm{d}x}{\mathrm{d}t}+\cos\theta\tfrac{\mathrm{d}y}{\mathrm{d}t}+\cos\varphi\sin\theta\tfrac{\mathrm{d}z}{\mathrm{d}t}\\
(\tfrac{\mathrm{d}Z}{\mathrm{d}x})\tfrac{\mathrm{d}x}{\mathrm{d}t}+(\tfrac{\mathrm{d}Z}{\mathrm{d}y})\tfrac{\mathrm{d}y}{\mathrm{d}t}+(\tfrac{\mathrm{d}Z}{\mathrm{d}z})\tfrac{\mathrm{d}z}{\mathrm{d}t}&=\zeta=-\cos\varphi\tfrac{\mathrm{d}x}{\mathrm{d}t}+\sin\varphi\tfrac{\mathrm{d}z}{\mathrm{d}t}
\end{aligned}
\]
The total relative speed is consequently $=\BareRoot(\xi\xi+\eta\eta+\zeta\zeta)=u$. As one readily sees by expanding the expressions above, this becomes $=\BareRoot(\tfrac{\mathrm{d}x^2}{\mathrm{d}t^2}+\tfrac{\mathrm{d}y^2}{\mathrm{d}t^2}+\tfrac{\mathrm{d}z^2}{\mathrm{d}t^2})$. Air resistance is proportional to its square, so we set it $=Muu$ and resolve it along the three directions above into $Mu\xi$, $Mu\eta$, $Mu\zeta$.

We regard the Earth here as a spheroid of revolution. The direction of gravity therefore passes through the Earth's axis. Let the point where it meets that axis lie a distance $q$ above $C$, so that for that point $Z=q$.

If we further set the strength of gravitation $=p$ and put
\[
XX+YY+(Z-q)^2=rr
\]
then the principles of dynamics give
\SourceLeaf{506}\[
\left.\begin{aligned}
0&=\frac{\mathrm{dd}X}{\mathrm{d}t^2}+\frac{pX}{r}+Mu\xi\\
0&=\frac{\mathrm{dd}Y}{\mathrm{d}t^2}+\frac{pY}{r}+Mu\eta\\
0&=\frac{\mathrm{dd}Z}{\mathrm{d}t^2}+\frac{p(Z-q)}{r}+Mu\zeta
\end{aligned}\right\}\quad.\quad.\quad.\quad.\quad.\quad.\quad.\quad[3]
\]
From the above values of $X,Y,Z$ in [2], the following equations result when the constant $\tfrac{\mathrm{d}\theta}{\mathrm{d}t}$ is written as $n$.
\[
\begin{aligned}
\frac{\mathrm{dd}X}{\mathrm{d}t^2}={}&\sin\varphi\cos\theta\frac{\mathrm{dd}x}{\mathrm{d}t^2}-\sin\theta\frac{\mathrm{dd}y}{\mathrm{d}t^2}+\cos\varphi\cos\theta\frac{\mathrm{dd}z}{\mathrm{d}t^2}\\
&-2n\sin\varphi\sin\theta\frac{\mathrm{d}x}{\mathrm{d}t}-2n\cos\theta\frac{\mathrm{d}y}{\mathrm{d}t}-2n\cos\varphi\sin\theta\frac{\mathrm{d}z}{\mathrm{d}t}-nnX\\[5pt]
\frac{\mathrm{dd}Y}{\mathrm{d}t^2}={}&\sin\varphi\sin\theta\frac{\mathrm{dd}x}{\mathrm{d}t^2}+\cos\theta\frac{\mathrm{dd}y}{\mathrm{d}t^2}+\cos\varphi\sin\theta\frac{\mathrm{dd}z}{\mathrm{d}t^2}\\
&+2n\sin\varphi\cos\theta\frac{\mathrm{d}x}{\mathrm{d}t}-2n\sin\theta\frac{\mathrm{d}y}{\mathrm{d}t}+2n\cos\varphi\cos\theta\frac{\mathrm{d}z}{\mathrm{d}t}-nnY\\[5pt]
\frac{\mathrm{dd}Z}{\mathrm{d}t^2}={}&-\cos\varphi\frac{\mathrm{dd}x}{\mathrm{d}t^2}+\sin\varphi\frac{\mathrm{dd}z}{\mathrm{d}t^2}
\end{aligned}\quad[4]
\]
First multiply the three equations [3] by $\sin\varphi\cos\theta$, $\sin\varphi\sin\theta$, $-\cos\varphi$, respectively, and add the products. Then multiply the same equations by $-\sin\theta$, $\cos\theta$, $0$, and in a third step by $\cos\varphi\cos\theta$, $\cos\varphi\sin\theta$, $\sin\varphi$, adding the products in each case. After substituting for $\tfrac{\mathrm{dd}X}{\mathrm{d}t^2}$, $\tfrac{\mathrm{dd}Y}{\mathrm{d}t^2}$, $\tfrac{\mathrm{dd}Z}{\mathrm{d}t^2}$ their values from [4], for $X,Y,Z$ their values from [2], and for $\xi$, $\eta$, $\zeta$ their respective expressions, one obtains the following three new equations.
\[
\begin{aligned}
0&=\tfrac{\mathrm{dd}x}{\mathrm{d}t^2}-2n\sin\varphi\tfrac{\mathrm{d}y}{\mathrm{d}t}+(x-a)(\tfrac pr-nn)+\cos\varphi(\tfrac{pq}{r}-nnZ)+Mu\tfrac{\mathrm{d}x}{\mathrm{d}t}\\
0&=\tfrac{\mathrm{dd}y}{\mathrm{d}t^2}+2n\sin\varphi\tfrac{\mathrm{d}x}{\mathrm{d}t}+2n\cos\varphi\tfrac{\mathrm{d}z}{\mathrm{d}t}+y(\tfrac pr-nn)+Mu\tfrac{\mathrm{d}y}{\mathrm{d}t}\\
0&=\tfrac{\mathrm{dd}z}{\mathrm{d}t^2}-2n\cos\varphi\tfrac{\mathrm{d}y}{\mathrm{d}t}+(z+c)(\tfrac pr-nn)-\sin\varphi(\tfrac{pq}{r}-nnZ)+Mu\tfrac{\mathrm{d}z}{\mathrm{d}t}
\end{aligned}
\]
If the body is thus at rest relative to the Earth, or $\mathrm{d}x=\mathrm{d}y=\mathrm{d}z=0$, it appears to be acted on, perpendicularly to the planes for $x,y,z$, by the forces
\SourceLeaf{507}\[
\begin{gathered}
(x-a)(\tfrac pr-nn)+\cos\varphi(\tfrac{pq}{r}-nnZ)\\
y(\tfrac pr-nn)\\
(z+c)(\tfrac pr-nn)-\sin\varphi(\tfrac{pq}{r}-nnZ)
\end{gathered}
\]
A body already in motion, however, is affected differently. Air resistance drives it along these directions like forces of magnitude $Mu\tfrac{\mathrm{d}x}{\mathrm{d}t}$, $Mu\tfrac{\mathrm{d}y}{\mathrm{d}t}$, $Mu\tfrac{\mathrm{d}z}{\mathrm{d}t}$, and therefore acts on the rotating Earth exactly as it would on an Earth at rest. In addition to that resistance, the three further forces
\[
-2n\sin\varphi\frac{\mathrm{d}y}{\mathrm{d}t},\quad 2n\sin\varphi\frac{\mathrm{d}x}{\mathrm{d}t}+2n\cos\varphi\frac{\mathrm{d}z}{\mathrm{d}t},\quad -2n\cos\varphi\frac{\mathrm{d}y}{\mathrm{d}t}
\]
act in those directions, and it is these alone that make the Earth's rotation visible in falling bodies. The reasoning and conclusions so far are strictly and generally valid.

In \textit{experiments} made with this in view, the body's motion always takes place within so small a region that one can take the strength of the apparent gravity acting on bodies at rest there to be constant $=g$, and its direction to be everywhere parallel, and hence perpendicular to the plane for $z$. One may therefore freely replace the three quantities above,
\[
\begin{gathered}
(x-a)(\tfrac pr-nn)+\cos\varphi(\tfrac{pq}{r}-nnZ)\\
y(\tfrac pr-nn)\\
(z+c)(\tfrac pr-nn)-\sin\varphi(\tfrac{pq}{r}-nnZ)
\end{gathered}
\]
by $0,0,g$, respectively. The three fundamental equations then become
\[
\begin{aligned}
0&=\frac{\mathrm{dd}x}{\mathrm{d}t^2}-2n\sin\varphi\frac{\mathrm{d}y}{\mathrm{d}t}+Mu\frac{\mathrm{d}x}{\mathrm{d}t}\\
0&=\frac{\mathrm{dd}y}{\mathrm{d}t^2}+2n\sin\varphi\frac{\mathrm{d}x}{\mathrm{d}t}+2n\cos\varphi\frac{\mathrm{d}z}{\mathrm{d}t}+Mu\frac{\mathrm{d}y}{\mathrm{d}t}\\
0&=\frac{\mathrm{dd}z}{\mathrm{d}t^2}-2n\cos\varphi\frac{\mathrm{d}y}{\mathrm{d}t}+g+Mu\frac{\mathrm{d}z}{\mathrm{d}t}
\end{aligned}
\]
\SourceLeaf{508}These equations are easy to integrate if air resistance is neglected, or $M=0$ is set. One finds
\[
\begin{aligned}
x&=\mathfrak{A}-\mathfrak{D}\cos\varphi.t+\tfrac{1}{2n}\mathfrak{E}\sin\varphi\cos(2nt+\mathfrak{F})+\tfrac12\sin\varphi\cos\varphi.gtt\\
y&=\mathfrak{B}-\tfrac{1}{2n}\mathfrak{E}\sin(2nt+\mathfrak{F})+\tfrac{1}{2n}\cos\varphi.gt\\
z&=\mathfrak{C}+\mathfrak{D}\sin\varphi.t+\tfrac{1}{2n}\mathfrak{E}\cos\varphi\cos(2nt+\mathfrak{F})-\tfrac12\sin\varphi^2.gtt
\end{aligned}
\]
It is also easy to derive the following values of the arbitrary quantities if we assume that the body initially has no apparent velocity at all.
\[
\begin{array}{lll}
\mathfrak{A}=-\tfrac{g}{4nn}\cos\varphi\sin\varphi,&\mathfrak{B}=0,&\mathfrak{C}=-\tfrac{g}{4nn}\cos\varphi^2\\
\mathfrak{D}=0,&\mathfrak{E}=\tfrac{g}{2n}\cos\varphi,&\mathfrak{F}=0
\end{array}
\]
Thus
\[
\begin{aligned}
x&=\tfrac{g}{2n}\cos\varphi\sin\varphi(ntt-\tfrac{1}{2n}+\tfrac{1}{2n}\cos2nt)\\
y&=\tfrac{g}{2n}\cos\varphi(t-\tfrac{1}{2n}\sin2nt)\\
z&=-\tfrac12gtt+\tfrac{g}{2n}\cos\varphi^2(ntt-\tfrac{1}{2n}+\tfrac{1}{2n}\cos2nt)
\end{aligned}
\]
This integration is of course not \textit{generally} permissible, since the above assumption is allowed only so long as the body does not move far from its initial apparent position. In this case, however, we can expand the trigonometric functions in series, so that
\[
\begin{aligned}
x&=\tfrac16\cos\varphi\sin\varphi.gnnt^4\quad.\quad.\quad.\\
y&=\tfrac13\cos\varphi.gnt^3\quad.\quad.\quad.\\
z&=-\tfrac12gtt+\tfrac16\cos\varphi^2.gnnt^4\quad.\quad.\quad.
\end{aligned}
\]
Since the fall lasts only a few seconds, and $nt$ therefore amounts to at most a few minutes of arc, and since (with radius $=1$) $1'=\tfrac{1}{3438}$, $x$ and the second part of $z$ are quite imperceptible, so $y=-\tfrac23z\cos\varphi.nt$. In Dr.\ \textsc{Benzenberg}'s experiment in the tower of St. Michael's, $z=-235$ feet, $\varphi=53^0\,33'$, and $t=4''$ of solar time, so $nt=\tfrac{366}{365}$ minutes of arc. This gives $y=3{,}91$ lines.

\SourceLeaf{509}If we want to take air resistance into account when integrating the above equations, we shall have to be content with approximations. It is then very easy to expand the values of $x,y,z$ in series in powers of $n$ and $M$. The highest term of $x$ is again $=\tfrac16\cos\varphi\sin\varphi.gnnt^4$, as before, and is therefore of no significance at all. For $y$ and $z$, neglecting the squares and higher powers of $n$ and $M$, one finds the following values.
\[
\begin{aligned}
y&=\tfrac13\cos\varphi.gnt^3-\tfrac{1}{12}\cos\varphi.Mggnt^5\\
z&=-\tfrac12gtt+\tfrac{1}{12}Mggt^4.
\end{aligned}
\]
Let us therefore put $-z$, the actual distance fallen, $=f$, and $\tfrac12gtt$, or the distance fallen in a vacuum, $=f+\delta$. Then
\[
y=\tfrac23\cos\varphi.nt(f+\delta)-\cos\varphi.nt\delta=\tfrac23\cos\varphi.nt(f-\tfrac12\delta)
\]
For the experiments at St.\ Michael's, where $f+\delta=241{,}47$ feet, we thus obtain the eastward deflection $y=3{,}86$ lines.

